% !TeX root = ../se200_identity_regimes.tex

\section{Conclusion}
\label{sec:conclusion}

\section{Closing Synthesis}

This paper establishes that admissible substrate systems admit
a neutral, regime-typed representation in which structure is realized
as families of admissible paths.

The construction is:

\begin{itemize}
    \item structurally sufficient for admissible multi-stage systems,
    \item faithful up to regime-typed isomorphism,
    \item neutral with respect to causal and normative interpretation.
\end{itemize}

Interpretation is deferred to extensions over REC-anchored path families.
Comparison reduces to path comparison.
Intervention reduces to admissible path transformation.

The representation does not encode causal or normative force.
Where such commitments are needed, they are introduced by extension layers
rather than by the substrate itself.

The result establishes that at least nine pairwise non-collapsing
regime profiles are required under the stated assumptions, and that
admissible structure within the finite-system setting can be represented
as regime-typed path families.
This provides a representational basis for subsequent analysis over
path structures subject to the admissibility constraints adopted here.

This paper derives structural constraints on identity and persistence
for neutral substrates under persistent interpretive disagreement.
Within this setting, identity is determined by invariance under admissible
transformations rather than by causal explanation or normative evaluation.

We show that transformation-induced split pressure requires refinement
of certain regime profiles.
Enduring non-normative referents split under branching,
applicability contexts split under decomposition,
and normative structures split under refinement,
each producing distinct profiles required for stable identity.


These results establish a necessary lower bound on the number of
referential regimes required for such substrates.
At least nine distinct regime profiles are required.
Completeness or minimality of this set is not established.

The result is structural and conditional.
It characterizes constraints that any substrate must satisfy
in order to support stable reference under admissible transformations,
without embedding causal or normative commitments.


% \section*{Acknowledgements}

% % The author thanks the anonymous reviewers for their detailed and
% % technical engagement with the argument.
% % Their comments substantially strengthened the necessity proof
% % and clarified the scope of the central result.

% Portions of this work were developed
% using computer-assisted tools during manuscript preparation.
% Generative language models were used to assist with editing, formatting,
% and consistency checking.
% All conceptual framing, formal development, results, interpretations, and conclusions
% are the author's own.

% The author reviewed all suggestions and
% takes full responsibility for the content of this work.
